% =====================================================================
\chapter{Le mode de rendu cinématique (VFX)}
\label{ch:vfx}
% =====================================================================

Ce chapitre décrit une fonctionnalité volontairement contradictoire avec le
reste du projet, et la manière dont la contradiction a été encadrée.

\section{Le problème : la tension}

Le produit est fondé sur l'honnêteté de la donnée : pas de score, pas
d'invention, une provenance sur chaque fait, une étiquette de résolution. Une
démonstration cinématique — caméra qui vole, eau qui monte avec des reflets,
particules, effets — est à l'opposé : elle vise l'impression, pas la preuve.

La tentation était double :

\begin{itemize}[leftmargin=1.4em]
  \item \textbf{commerciale} — une démonstration spectaculaire vend un concept
        mieux qu'un JSON ;
  \item \textbf{technique} — un rendu 3D temps réel est un objet d'ingénierie
        intéressant.
\end{itemize}

Le risque était tout aussi double : introduire un rendu qui \emph{paraît}
factuel alors qu'il est illustratif, et laisser ce rendu contaminer le contrat
de données.

\section{La solution : gouverner la contradiction}

Plutôt que de refuser la fonctionnalité ou de l'accepter telle quelle, le
projet l'a encadrée par un amendement constitutionnel. La constitution a été
portée en version~0.2 le 14 septembre 2026, avec un paragraphe §2.1 dédié.

\subsection{Deux modes obligatoires}

\begin{table}[H]
\centering
\small
\caption{Les deux modes de rendu et leurs contrats respectifs.}
\label{tab:modes}
\begin{tabularx}{\textwidth}{@{}l X X@{}}
\toprule
& \textbf{Mode \code{data}} & \textbf{Mode \code{vfx}} \\
\midrule
Rendu & Enveloppes honnêtes, extrusion Mapbox sur les empreintes réelles &
  Couche personnalisée Three.js sur Mapbox, terrain, particules,
  post-traitement \\
\addlinespace
Défaut & \textbf{oui} & non — activé explicitement \\
\addlinespace
Étiquetage & Aucun (le rendu est factuel) &
  Bandeau permanent \og Simulation visuelle — non contractuelle \fg \\
\addlinespace
Données partielles & Mention d'indisponibilité explicite &
  Étiquette secondaire \og Données partielles — rendu illustratif \fg \\
\addlinespace
Contrat API & Inchangé & \textbf{inchangé} \\
\bottomrule
\end{tabularx}
\end{table}

\begin{invariant}
Aucune route backend, aucun champ de schéma, aucune charge de lot ne peut
transporter de géométrie VFX, de score ou de jugement synthétique. Le mode VFX
est strictement confiné à la présentation.
\end{invariant}

\subsection{Règles d'entrée}

Le mode VFX doit consommer \emph{les mêmes} pilotes que le mode données
lorsqu'ils sont disponibles : l'heure simulée, les profils météorologiques, les
bandes de profondeur TRI, le trajet hydrographique, la hauteur et l'emprise
BDNB. Il peut exagérer l'\emph{apparence} — montée lissée, plan d'eau suivant
le terrain, floraison lumineuse, mouvement de caméra — mais il ne peut pas
revendiquer une précision hydraulique.

Les replis illustratifs sont autorisés \emph{uniquement} en mode VFX et
\emph{uniquement} avec l'étiquette de données partielles.

\subsection{Le bus d'entrée unique}

Un module dédié, \fichier{vfxInputs.ts}, est le seul point d'entrée du rendu
cinématique. La règle est explicite : le rendu VFX passe par ce module et rien
d'autre.

\begin{lstlisting}[caption={Le bus d'entrée du mode VFX et sa règle de repli.},
label={lst:vfxinputs}]
vfxWaterLevel({ rainProfile, hourIndex, triPeak, allowFallback })
  -> { levelM, partial, source: 'forecast' | 'fallback' | 'none' }
\end{lstlisting}

Deux garanties ont été codées dans ce module :

\begin{enumerate}[leftmargin=1.6em]
  \item Le repli illustratif est \textbf{plafonné par la classe TRI} :
        $\min(\text{triPeak}, 1{,}0~\text{m})$. Le rendu ne peut donc jamais
        exagérer au-delà de l'exposition réglementaire réelle.
  \item \textbf{Sans classe TRI, la source vaut \code{none} et le niveau
        vaut 0}, même si le repli est autorisé. On n'invente pas une exposition
        qui n'existe pas.
\end{enumerate}

Ce point est important : le mode cinématique peut embellir la \emph{forme},
jamais créer un \emph{fait}. Une adresse hors zone inondable reste sèche, même
en mode spectaculaire.

\begin{honesty}
Un problème d'incohérence a été détecté et corrigé pendant la conception : la
spécification initiale prévoyait une constante de repli de 1,0~m indépendante
de la classe TRI. Elle était à la fois incohérente — elle pouvait afficher plus
que le risque réel — et inatteignable, puisque le mode données consommait une
autre fonction. Elle a été supprimée : le repli ne porte plus que sur la
\emph{forme} de la courbe.
\end{honesty}

\section{Contraintes d'intégration}

\subsection{Exclusion mutuelle avec la couche météo}

La couche météo Windy impose une projection plate (Mercator sans inclinaison).
La caméra 3D du mode VFX, elle, exige une inclinaison forte. Les deux sont
incompatibles.

\begin{invariant}
La couche Windy et la caméra 3D du mode VFX sont mutuellement exclusives tant
que le mode VFX est actif. Cette exclusivité est vérifiée par un test.
\end{invariant}

Concrètement : à l'activation du mode VFX, la couche Windy est désactivée et
l'inclinaison de la carte est forcée au-delà de 55 degrés. À la sortie, la
politique du mode données reprend la main et l'état précédent est restauré.

\subsection{Le montage et le démontage}

Le mode VFX doit s'activer et se désactiver proprement. La contrainte est
double : aucune fuite de contexte WebGL au changement de route, et aucun
doublon de couche en cas de bascule répétée.

\begin{invariant}
La couche VFX se monte et se démonte proprement lors de la bascule de mode.
Aucune fuite de contexte WebGL au changement de route. Objectif de performance :
au moins 30~images par seconde sur un ordinateur portable de milieu de gamme au
niveau de zoom d'un bâtiment avec le plan d'eau actif, avec dégradation
gracieuse — le post-traitement est désactivé en premier.
\end{invariant}

\section{Étiquetage}

L'étiquetage n'est pas cosmétique : c'est la condition qui rend la
fonctionnalité acceptable.

\begin{itemize}[leftmargin=1.4em]
  \item Le bandeau \og \textbf{Simulation visuelle — non contractuelle} \fg{}
        est affiché en permanence en mode VFX.
  \item Une étiquette secondaire \og \textbf{Données partielles — rendu
        illustratif} \fg{} apparaît lorsque les entrées réelles sont
        incomplètes.
  \item Le mode VFX \textbf{ne doit pas réutiliser le vocabulaire du contrat}
        (\code{resolution}, \code{present}, libellés de bandes) comme s'il
        s'agissait d'un fait sourcé.
\end{itemize}

Le composant d'avertissement ne s'affiche \emph{que} en mode VFX : il n'existe
pas de bandeau dormant en mode données.

\section{État d'avancement}

Toutes les tâches du mode VFX n'ont pas été menées à terme pendant la durée du
stage. Le tableau~\ref{tab:vfx-etat} distingue ce qui est livré de ce qui
reste, sans embellissement.

\begin{table}[H]
\centering
\small
\caption{État d'avancement des tâches du mode VFX.}
\label{tab:vfx-etat}
\begin{tabularx}{\textwidth}{@{}l l X@{}}
\toprule
\textbf{Tâche} & \textbf{État} & \textbf{Contenu} \\
\midrule
\vfx{001} & Livré & Amendement constitutionnel, spécification et plan \\
\vfx{002} & Livré & Bascule de mode, bandeau, mise en page atténuée \\
\vfx{003} & Livré & Bus d'entrée, pilote de niveau d'eau, plafonnement TRI \\
\vfx{006} & Livré & Politique de caméra, exclusion mutuelle Windy \\
\vfx{008} & Livré (partiel) & Isolation des modes, pas de double nappe d'eau \\
\vfx{004} & Restant & Dépendance Three.js et couche personnalisée Mapbox \\
\vfx{005} & Restant & Grille de terrain et nuanceur d'eau \\
\vfx{007} & Restant & Réalisateur caméra (\og Play scenario \fg{}) \\
\vfx{009} & Optionnel & Post-traitement \\
\vfx{010} & Optionnel & Effets secondaires vent et feu \\
\bottomrule
\end{tabularx}
\end{table}

Un point d'honnêteté sur cet état : le niveau d'eau calculé par le mode VFX
n'alimente pas encore de couche de rendu 3D, puisque la tâche \vfx{004} n'est
pas achevée. En revanche, ce niveau n'est pas du code dormant : il alimente
l'étiquette de données partielles du bandeau, ce qui est un usage réel — la
constitution interdit le code non référencé.

\section{Bilan de la démarche}

L'intérêt de ce chapitre n'est pas le rendu cinématique lui-même, mais la
manière dont une tension produit a été traitée. Trois éléments ont rendu la
fonctionnalité acceptable :

\begin{enumerate}[leftmargin=1.6em]
  \item \textbf{un amendement explicite} — la règle a été changée avant le
        code, pas après ;
  \item \textbf{une étanchéité vérifiée par des tests} — le mode VFX ne peut
        pas modifier le contrat ;
  \item \textbf{un étiquetage obligatoire} — le spectateur sait toujours qu'il
        regarde une illustration.
\end{enumerate}

Sans ces trois éléments, la fonctionnalité aurait été soit refusée, soit
introduite en silence — et dans le second cas, elle aurait probablement érodé
la confiance dans l'ensemble du produit.
